%%%%%%%%%%%%%%%%%%%%%%%%%%%%%%%%%
%
%       EXERCÍCIO
%
%%%%%%%%%%%%%%%%%%%%%%%%%%%%%%%%%

\definevalorquestao

\begin{exercicioBanco}[\valorquestao]
Um concorrente participa de um jogo de auditório no qual há três portas fechadas ($A$, $B$ e $C$). Atrás de uma das portas há um carro zero quilômetro e, atrás das outras duas, há apenas uma ovelha em cada. O prêmio do participante é o que estiver atrás da porta escolhida. 

O jogo ocorre da seguinte maneira:
\begin{enumerate}[1.]
    \item O participante escolhe inicialmente uma porta (por exemplo, a Porta $A$).
    \item O apresentador, que sabe exatamente onde está o carro, abre uma das outras duas portas que ele sabe que contém uma ovelha (por exemplo, a Porta $C$).
    \item Em seguida, o apresentador faz a seguinte pergunta ao participante: \textit{“Você deseja manter a sua escolha original ou prefere trocar para a porta que sobrou (Porta $B$)?”}
\end{enumerate}

Com base nas regras de probabilidade, responda aos itens a seguir:

\begin{enumerate}[(a)]
    \item Qual é a probabilidade de o participante ganhar o carro se ele não trocar de porta?
    \item Qual é a probabilidade de ele ganhar o carro se ele aceitar trocar de porta? Justifique calculando as probabilidades condicionais associadas.
    \item A estratégia de trocar de porta realmente aumenta as chances do concorrente de ganhar o prêmio?
\end{enumerate}
\end{exercicioBanco}


